% Supplementary Information of the ERL Letter (D-107 addendum): how the main display items are
% drawn. These details stood in the Letter's legends until the legends were shortened to at most
% 120 words at the PI's request; they come from figures/erl_fig_*.py, figures/_erl.py and
% scripts/f1_erl_display.py (D-102). Every number is a \metric{} macro.

\subsection*{Display items of the Letter}
\emph{Maps.} Every map uses the Lambert azimuthal equal-area projection centred at \metric{params.laea_lat0_deg}\textdegree{}N, \metric{params.laea_lon0_deg}\textdegree{}E (Great Salt Lake: centred on the lake), with graticules every \metric{params.erl_fig_graticule_map_deg}\textdegree{} on whole-domain maps and \metric{params.erl_fig_graticule_zoom_deg}\textdegree{} on zooms and on Great Salt Lake, and scale bars of \metric{params.erl_fig_scale_map_km}~km (domain), \metric{params.erl_fig_scale_inset_km}~km (zooms) and \metric{params.erl_fig_scale_gsl_km}~km (Great Salt Lake). Light blue marks the open sea of the \metric{params.lowstand_label} low stand from Landsat, for orientation. The Supplementary maps use the same projection, with a graticule every \metric{params.fig_graticule_deg}\textdegree{} and scale bars of \metric{params.fig_scale_bar_km}~km, true at the projection centre.

\emph{Then and now (figure~\ref{M-fig:erl_thennow} of the Letter).} The composites are true-colour medians of Landsat Collection~2 Level-2 surface reflectance\citep{usgs_l45tm_c2l2,usgs_l89oli_c2l2} at \metric{erl.landsat.res_m}~m, with clouds, cloud shadows and saturated pixels removed and one reflectance stretch for both, so their colours compare. Pixels without a clear value in the main summer are taken from the summer before: \metric{erl.landsat.tm_1996.filled_pct}\% from \metric{erl.landsat.tm_1996.fill_year} ($n = \metric{erl.landsat.tm_1996.n_scenes_fill}$ scenes) for \metric{erl.landsat.tm_1996.year}, and \metric{erl.landsat.oli_2025.filled_pct}\% from \metric{erl.landsat.oli_2025.fill_year} for \metric{erl.landsat.oli_2025.year}. The shorelines are contours of the emergence date (section ``Waterlines and seabed topography'') of the seabed between the \metric{params.highstand_label} and \metric{params.lowstand_label} stands; for display the dates are block medians over \metric{params.erl_fig_contour_block}~$\times$~\metric{params.erl_fig_contour_block} pixels of \metric{params.fig_map_res_m}~m, lightly smoothed, and contour pieces shorter than \metric{params.erl_fig_contour_min_km}~km are omitted. Circles on the level record mark the level at the start of each contour year. The level record runs from \metric{levels.dahiti.first} to \metric{levels.dahiti.last}, with Hydroweb levels shifted onto the DAHITI reference in its gaps.

\emph{Shoreline map (figure~\ref{M-fig:erl_shorelines}).} The observed state (sea connected to the open sea, land, or not determined: too few valid observations, or water not connected to the sea that touches undetermined cells) is drawn on a \metric{erl.hero.display_res_m}~m grid as the majority of its \metric{params.stage2_res_m}~m cells; the zooms are at \metric{params.stage2_res_m}~m. The predicted shorelines are those of the shape variant, with constants of \metric{h1.shape_offset.gebco2026_m}~m (GEBCO\_2026), \metric{h1.shape_offset.etopo2022_m}~m (ETOPO 2022) and \metric{h1.shape_offset.srtm15p_m}~m (SRTM15+); where lines coincide, the last drawn hides the others. The two levels are bins \metric{params.bin_width_m}~m wide of E3, one on each side of \metric{h1.disp_band.edge_2_m}~m. Each arrow runs, for one grid, from the observed shoreline cell of the zoom's sector whose distance to the grid's nearest predicted shoreline cell is closest to that sector's median displacement at that level, to that cell. The Volga zoom at the higher level holds no such cell near the sector median, so its arrows are longer than the median. All $n = \metric{erl.hero.n_checks}$ sector medians of the two levels were recomputed from the observations and equal the audit's.

\emph{Audit figure (figure~\ref{M-fig:erl_audit}).} Panel (a) counts every observed level of every epoch inside the domain (\metric{h1.n_units_domain.E1} in E1, \metric{h1.n_units_domain.E2} in E2 and \metric{h1.n_units_domain.E3} in E3), as a percentage of the observed sea area in the cells observed at that level. Panel (b) uses the band edges \metric{h1.disp_band.edge_0_m}, \metric{h1.disp_band.edge_1_m}, \metric{h1.disp_band.edge_2_m}, \metric{h1.disp_band.edge_3_m}, \metric{h1.disp_band.edge_4_m} and \metric{h1.disp_band.edge_5_m}~m, and the median over sectors and levels of the per-level median displacement. The per-block maps of every grid are figures~\ref{fig:ed_h1blocks_shape} and \ref{fig:ed_h1blocks_total}, where blocks with fewer than \metric{params.h1_min_block_n} distances are not drawn. Panel (c) gives, per sector, the median over evaluable levels of each grid's per-level median displacement, against the median over the same levels of the noise-floor percentile.

\emph{Great Salt Lake (figure~\ref{M-fig:erl_gsl}).} The level drawn is the observation unit, among the $n = \metric{erl.gsl.n_candidate_units}$ observed in both arms, whose displacements lie closest to the summary medians of panel (b); a first rule, which picked an atypical unit, was replaced before the figure was drawn (section~\SIdeviations). At that unit the median displacement is \metric{erl.gsl.north.gebco2026.median_m}~m (north arm) and \metric{erl.gsl.south.gebco2026.median_m}~m (south arm) for GEBCO\_2026, \metric{erl.gsl.north.etopo2022.median_m} and \metric{erl.gsl.south.etopo2022.median_m}~m for ETOPO 2022, \metric{erl.gsl.north.srtm15p.median_m} and \metric{erl.gsl.south.srtm15p.median_m}~m for SRTM15+, and \metric{erl.gsl.north.dep3.median_m} and \metric{erl.gsl.south.dep3.median_m}~m for the lidar. The lidar is evaluated at \metric{generality.great_salt_lake.control.north.n_evaluable_levels} and \metric{generality.great_salt_lake.control.south.n_evaluable_levels} levels above its flattened lake surface only, with noise floors of \metric{erl.gsl.control.north.noise_floor_m} and \metric{erl.gsl.control.south.noise_floor_m}~m at those levels; the grids' floors are \metric{erl.gsl.north.noise_floor_m}~m (north) and \metric{erl.gsl.south.noise_floor_m}~m (south).

\emph{Agreement map (figure~\ref{M-fig:erl_agreement}).} The count of equation~(\ref{eq:si_agreement}) is drawn on a \metric{erl.hero.display_res_m}~m display grid as the most frequent count among its cells. At least three grids place \metric{erl.agree.1m.area_km2_majority}~\kmsq{} in the band at \metric{h3.scenario.1_m}~m and \metric{erl.agree.2m.area_km2_majority}~\kmsq{} at \metric{h3.scenario.2_m}~m. The symbols mark OpenStreetMap features (data \textcopyright{} OpenStreetMap contributors, Open Database License)\citep{osm} within \metric{params.erl_fig_near_km}~km of a cell that any grid places in the band at \metric{h3.scenario.2_m}~m.

\emph{Report card (table~\ref{M-tab:reportcard}).} The shading cut-offs are fixed: sectors mispredicted, good at most \metric{params.erl_card_sectors_good} and poor from \metric{params.erl_card_sectors_poor}; displacements, good below \metric{params.erl_card_km_good}~km and poor from \metric{params.erl_card_km_poor}~km (the range by its largest sector); height deviation, good below \metric{params.erl_card_nmad_good_m}~m and poor from \metric{params.erl_card_nmad_poor_m}~m; Great Salt Lake arms, good \metric{params.erl_card_arms_good} and poor \metric{params.erl_card_arms_poor}. The displacement in the lowest level band is the median over sectors and levels below \metric{h1.disp_band.edge_1_m}~m, the seabed exposed last.

\emph{Timeline (figure~\ref{M-fig:erl_timeline}).} Each pathway is its own change after its base year, applied from the 2026 level and linear between the years its source states: the best-fit SSP2-4.5 and SSP5-8.5 totals and the observed 1996--2021 trend carried forward\citep{samant2023}, and the medians of the low, medium, medium-to-high and high emission scenarios\citep{hoseini2025}. The first metre is reached between \metric{erl.text.h3.paths.1m.first_year} and \metric{erl.text.h3.paths.1m.last_year}; the low-emission median reaches none of the scenario depths by 2100. Shading marks the first and second metre below the 2026 level.
